% Einheit 3 von 3 – Quadratwurzeln (bank/potenzen-wurzeln/e3.jsonl, zusammenbau v0.1)
%% TODO Zweigzeile Teil 1: was hier gelernt wird (Satz in Schülersprache aus dem Einheitstitel) – Einheit 3
\einheitenkopf[e3]{Einheit 3 von 3 · Quadratwurzeln}
\zweigzeile{ab Kl. 7, je nach Buch bis Kl. 9 (am Gymnasium ab Kl. 8) · P10}
\setcounter{aufgabe}{20}

%% TODO Titel: Ich-kann-Satz fehlt in der Bank (Platzhalter: Kettenname) – Nr. 21
%% TODO Anweisung: ein Satz über der ersten Teilaufgabe fehlt in der Bank – Nr. 21
\begin{aufgabe}{Was zuerst?}
\begin{teile}
\teil $\sqrt{(-7)^2}$ – was zuerst?\\ \kreuz{erst die Wurzel} \\ \kreuz{erst das Quadrat}
\end{teile}
\end{aufgabe}

%% TODO Titel: Ich-kann-Satz fehlt in der Bank (Platzhalter: Kettenname) – Nr. 22
%% TODO Anweisung: ein Satz über der ersten Teilaufgabe fehlt in der Bank – Nr. 22
\begin{aufgabe}{Quadratwurzeln}
%% TODO \rechenplatz{n}: Schrittzahl des Grundfalls fehlt in der Bank (nur bei fester Schreibform, 2.3 b)
\begin{teile}
\teil $\sqrt{121}$ – geht die Wurzel im Kopf auf?\\ \kreuz{ja, Quadratzahl} \\ \kreuz{nein, Näherungswert}
\teil $\sqrt{49}$ – wie viel? Mache die Probe. \leerfeld , Probe: \leerfeld
\teil $\sqrt{11}$ – mit dem Taschenrechner, auf zwei Stellen nach dem Komma gerundet? \leerfeld
\teil $\sqrt{60}$ – zwischen welchen ganzen Zahlen? Prüfe dann mit dem Taschenrechner. zwischen \leerfeld und \leerfeld
\teil $\sqrt{3} \;\square\; 1{,}7$ – $<$, $=$ oder $>$?
\teil $\sqrt{7}$; $2{,}6$; $\frac{5}{2}$; $-2{,}7$ – aufsteigend geordnet?
\teil $\left(\sqrt{13}\right)^2$ – wie viel? \leerfeld
\teil $\sqrt{(-8)^2}$ – wie viel? \leerfeld
\teil $\sqrt{-49}$ und $-\sqrt{49}$ – welcher Term hat einen Wert? Gib ihn an. \leerfeld
\teil $\sqrt{0{,}49}$ – wie viel? \leerfeld
\teil Ein quadratischer Teppich hat einen Flächeninhalt von $6{,}25\,\text{m}^2$. Wie lang ist eine Seite? \hfill (P10 2020 OS) \\ \leerfeld[m]
\teil $\sqrt{7^2 + 24^2}$ – wie viel? \leerfeld
\teil $\sqrt[3]{64}$ – wie viel? \leerfeld
\teil Welche Aussage ist wahr? Kreuze an. \hfill (P10 2015 OS)\\ \kreuz{$\sqrt{(-7)^2} = -7$} \\ \kreuz{$\sqrt{(-7)^2} = 7$} \\ \kreuz{$\sqrt{(-7)^2}$ ist nicht definiert}
\end{teile}
\end{aufgabe}

%% TODO Titel: Ich-kann-Satz fehlt in der Bank (Platzhalter: Kettenname) – Nr. 23
%% TODO Anweisung: ein Satz über der ersten Teilaufgabe fehlt in der Bank – Nr. 23
\begin{aufgabe}{Quadratwurzeln – Fehler finden, Begründen}
\begin{teile}
\teil Emil rechnet $\sqrt{64} = 32$. Wo ist der Fehler?
\teil Begründe, warum $\sqrt{49}$ nicht $-7$ ist, obwohl $(-7)^2 = 49$ gilt.
\end{teile}
\end{aufgabe}
